% Figure: a cylindrical (barrel) shell roof carrying its self-weight, showing
% the schematic geometry and the membrane stress-resultant sign along the arc.
% A single cross-section of a circular cylindrical shell of radius R spanning an
% angle from -phi0 to +phi0, supported along its two straight edges. We draw the
% arc, the vertical self-weight arrows, and annotate the crown (longitudinal
% compression) versus the free edges (tension), the reason a shell roof needs
% edge beams.
\begin{figure}[htbp]
\centering
\begin{tikzpicture}[>=Latex, scale=1.0]
% arc of the shell: radius 3, centred at (0,0), from 120 deg to 60 deg
\def\Rr{3}
\draw[engblue, very thick] (120:\Rr) arc[start angle=120, end angle=60, radius=\Rr];
% centre of curvature
\fill[enggray] (0,0) circle (1pt);
\node[enggray, font=\scriptsize, anchor=north] at (0,-0.05) {centre of curvature};
% radius lines to the two edges and the crown
\draw[enggray!60, dashed] (0,0) -- (120:\Rr);
\draw[enggray!60, dashed] (0,0) -- (90:\Rr);
\draw[enggray!60, dashed] (0,0) -- (60:\Rr);
% angle label
\node[enggray, font=\scriptsize] at (105:1.05) {$\phi_0$};
% self-weight arrows hanging from the arc
\foreach \a in {66,78,90,102,114} {
  \draw[engred, ->] (\a:\Rr) -- ++(0,-0.55);
}
\node[engred, font=\scriptsize, anchor=west] at (60:\Rr) {self-weight $g$};
% crown point and label
\fill[engamber] (90:\Rr) circle (1.6pt);
\node[engamber, font=\scriptsize, anchor=south] at (90:\Rr) {crown: longitudinal compression};
% free edges
\fill[englime!60!black] (120:\Rr) circle (1.6pt);
\fill[englime!60!black] (60:\Rr) circle (1.6pt);
\node[englime!60!black, font=\scriptsize, anchor=east] at (120:\Rr) {free edge:};
\node[englime!60!black, font=\scriptsize, anchor=east] at ($(120:\Rr)+(0,-0.32)$) {tension};
% edge beam boxes under the two edges
\draw[engblack, thick, fill=enggray!15]
  ($(120:\Rr)+(-0.25,-0.05)$) rectangle ($(120:\Rr)+(0.25,-0.45)$);
\draw[engblack, thick, fill=enggray!15]
  ($(60:\Rr)+(-0.25,-0.05)$) rectangle ($(60:\Rr)+(0.25,-0.45)$);
\node[engblack, font=\scriptsize, anchor=north] at ($(60:\Rr)+(0,-0.5)$) {edge beam};
\end{tikzpicture}
\caption[Barrel-vault shell roof]{A circular cylindrical (barrel) shell roof
spanning a half-angle $\phi_0$ each side of the crown and supported along its
two straight edges. Under self-weight the shell acts as an arch across its
curvature and as a beam along its length. The membrane theory gives a
longitudinal compression that is largest at the crown (amber) and swings into
tension near the free straight edges (green), where a shell with no edge beam
would pull inward and the membrane state breaks down. Real barrel-vault roofs
carry edge beams along the free edges to absorb that edge tension and to keep
the shell in the membrane regime over its interior, the shell counterpart of
the wind girder that keeps a storage-tank shell round.}
\label{fig:vol03ch09:barrel-vault}
\end{figure}
